\section{Scope, conventions, and principal results}
\label{cont:sec:introduction}

This article continues the research programme in \emph{Polylogarithms
and their Arithmetic Bridges} at the immutable ProveIt commit
\texttt{cc34f73596336f2466d9754cb0f3635bd2bedade}, dated
10 October 2026 \cite{cont:ProveIt}. The baseline includes both the
consolidated manuscript and six newly deposited research archives.
This distinction matters: several conjectures still labelled open in
the consolidated text already have proofs in those incoming archives.
The present contribution does not count those proofs as new work.
An accompanying inventory records the exact source and archive hashes.

The results concern three structures which recur throughout the
programme. Positive Mellin integrals control alternating harmonic sums;
gamma concentration describes a transition in reflected log-gamma
moments; integral distribution relations organize rational-grid special
values. A shorter fourth part reconciles the two recorded weight-seven
Gaussian candidates by an exact shuffle identity. The proofs are analytic
and algebraic, with one explicitly identified computer-assisted inequality
whose finite middle-interval step uses exact integer arithmetic. Other
computations supply reproducible checks and examples, and are identified
explicitly when they provide only numerical evidence.

\subsection{What is already known in the pinned material}

The consolidated manuscript already proves the mixed Gaussian
$S_4$ evaluation, signed first-omitted-term bounds for the harmonic
family, its first-pole transition, an all-orders expansion within that
transition, and a saddle expansion when the exponent-to-depth ratio
stays in a fixed positive compact interval. Its reflected log-gamma
theory includes fixed-reflected-exponent expansions, late coefficients,
and a sharp optimal-truncation remainder. Its distribution theory
includes arbitrary complex weights, conductor coordinates, exceptional
spectral jets, and a primitive-grid determinant up to a constant from
the choice of complex coordinates.

The six incoming archives additionally include proofs of the odd-weight
rank conjecture for the signed depth-two product module, extensions to
other levels, integral Smith calculations for that product module,
independent $S_4$ certificates, complement-collision classifications,
CM evaluations, and further zero-geometry results. These are distinct
from the integral distribution lattice studied here. A small apparent
coefficient relation found numerically at a higher weight is not made
a theorem by being stored in an incoming report. In particular, both
recorded $S_6$ formulas remain conjectural.

\subsection{The contributions proved here}

\begin{enumerate}[leftmargin=*,label=\arabic*.]
\item \textbf{A single harmonic saddle over the whole moving range.}
Theorem~\ref{us:thm:uniform} gives an all-orders \emph{relative}
expansion as $p\to\infty$, uniformly for all integers $h\ge1$ and
all $a>0$. It allows $p/h$ to tend to zero, to remain bounded, or to
grow without bound. The exact saddle incorporates the shift $a$.
A global gamma-phase upper bound makes this uniformity possible.
This resolves the request in the proportional-depth chapter for an
approximation which continues through the growing-ratio transition.

\item \textbf{An inverse transition with explicit corrections.}
Theorem~\ref{us:thm:inverse} determines the unique real outer exponent
at which the normalized harmonic sum reaches a prescribed value
$\theta\in(0,1)$. It gives a finite recursion to every order and
prints the first two correction polynomials. Its error is proved by
asymptotic brackets, without differentiating an unspecified remainder.

\item \textbf{A transition when both reflected moment exponents grow.}
Theorem~\ref{joint:thm:transition} treats
$(n+1)/(m+1)=\log m+O(1)$ with $m\to\infty$. The fixed-$m$
endpoint formula acquires a surviving factor
\[
 \exp\!\left[\left(\frac{\zeta(2)}{2\gamma}-\gamma\right)
       m\exp\!\left(-\frac{n+1}{m+1}\right)\right].
\]
Two corrections and an all-orders coefficient rule are proved.
Corollary~\ref{joint:cor:top-coefficients} gives exact identities for
two extremal coefficient families in every order.

\item \textbf{A unique saddle for every proportional moment.}
Theorem~\ref{gs:thm:monotonicity} proves that
\[
 x\longmapsto
 \frac{-\psi(x)\log\Gamma(1-x)}
      {-\psi(1-x)\log\Gamma(x)}
\]
is strictly increasing on $(0,1)$. The proof combines analytic
endpoint estimates with a small exact rational interval certificate.
Theorem~\ref{gs:thm:proportional} consequently gives the complete
moment expansion for exponent ratios in any compact positive
interval, including a unique nondegenerate saddle. This settles
another part of the simultaneous large-exponent problem.

\item \textbf{Integral coordinates and the complete two-prime
primitive cokernel.} Theorem~\ref{intdist:thm:basis} gives a concrete
grid basis over every commutative coefficient ring. Its finite
resolution is an explicit adaptation of classical universal-distribution
methods. The principal additional calculation is
Theorem~\ref{intdist:thm:two-prime}: a complete presentation of the
primitive-grid cokernel at $q=p\ell$, including torsion at singular
weights. Its Smith factors give the exact denominator required to
express all grid points in primitive coordinates. Explicit level-fifteen
Hurwitz-zeta identities illustrate the result.

\item \textbf{An exact reconciliation of the weight-seven
candidates.} Proposition~\ref{s6audit:prop:coordinate} proves a
coordinate identity among Gaussian doubles and shows that the two
$S_6$ candidates are the same conjecture. The displayed certificate
uses two convergent shuffles. An incomplete larger relation search
is reported as inconclusive.
\end{enumerate}

The phrase ``additional result'' means an extension of the audited
baseline, with the antecedents stated below. It is not a claim of
global priority after an exhaustive literature search. In particular,
the freeness/resolution principle in universal distributions is
classical; its role here is to support explicit integral calculations,
not to rename an existing general theory.

\subsection{Notation and interpretation}

We use $\gamma$ for Euler's constant and
$\beta(s)=\sum_{n\ge0}(-1)^n(2n+1)^{-s}$ for the Dirichlet beta
function when the defining series converges. The multiple-polylogarithm
convention is
\[
 \operatorname{Li}_{s_1,\ldots,s_d}(z_1,\ldots,z_d)
 =\sum_{n_1>\cdots>n_d\ge1}
    \frac{z_1^{n_1}\cdots z_d^{n_d}}
         {n_1^{s_1}\cdots n_d^{s_d}},
\]
continued only where a continuation is expressly indicated. All
Gaussian shuffles used for the coordinate theorem converge ordinarily.
The normalized harmonic quantity $R_{p,h}$ and the normalized moment
$\mathcal R_{n,m}$ are different functions. The symbol $\lambda$ is a
local transition parameter defined separately in each section.

The term \emph{formal distribution module} refers to symbols subject
to specified linear multiplication relations. Its rank is not the
dimension of the span of the corresponding evaluated special values.
Similarly, a Smith factor describes a particular integral inclusion;
it is not an assertion of algebraic or transcendence independence.
These distinctions are essential at exceptional weights and in finite
characteristic.

\subsection{Reading and integration guide}

Sections~\ref{us:sec:uniform}--\ref{us:sec:inverse} extend the harmonic
chapters. Sections~\ref{joint:sec:moments}--\ref{gs:sec:main} belong
after the fixed-$m$ reflected-moment and sharp-remainder material.
Section~\ref{intdist:sec:main} supplements the distribution and
conductor chapters. Section~\ref{s6audit:sec:main} belongs beside the
weight-seven conjecture. Section~\ref{cont:sec:editorial} records
proposed corrections and status changes; the final sections describe
the computational evidence and a prioritized research agenda.
The archive includes both an assembled \TeX{} article and modular
section files, together with exact verifiers, numerical diagnostics,
source receipts, and a small proposed editorial patch.
